\chapter{Alcance, autonomía y masa variable}\label{ch:breguet}
\objetivos{Derivar las integrales de consumo; asociar cada expresión de Breguet a una ley de vuelo; optimizar correctamente alcance y tiempo; actualizar una misión por fases.}

\section{La ecuación que debe preceder a cualquier Breguet}
En crucero horizontal cuasiestacionario, $T=D$ y
\begin{equation}
\dot m=-c_TD,\qquad \dot W=-g_0c_TD.
\end{equation}
Entre pesos inicial y final $W_i>W_f>0$, las integrales generales son
\begin{equation}
E=\int_{W_f}^{W_i}\frac{\dd W}{g_0c_TD},\qquad
R_a=\int_{W_f}^{W_i}\frac{V\,\dd W}{g_0c_TD}.
\label{eq:jet-general}
\end{equation}
En ellas $V$, $c_T$, $D$ y la densidad deben evaluarse según la ley de vuelo. El logaritmo característico no aparece por el nombre Breguet; aparece cuando el denominador es proporcional al peso y las demás magnitudes necesarias pueden sacarse de la integral.

\section{Caso I: velocidad y eficiencia constantes}
Si $E_a=L/D$ es constante, $D=W/E_a$. Si además $c_T$ y $V$ son constantes,
\begin{equation}
E=\frac{E_a}{g_0c_T}\ln\frac{W_i}{W_f},\qquad
R_a=\frac{VE_a}{g_0c_T}\ln\frac{W_i}{W_f}.
\label{eq:breguet-log}
\end{equation}
La eficiencia constante se obtiene con $C_L$ constante si la polar permanece igual. Con $V$ y $C_L$ constantes, la condición $L=W$ exige $\rho\propto W$. Al consumirse combustible, la densidad debe bajar: aparece un crucero con ascenso progresivo, siempre que la atmósfera y la envolvente permitan seguirlo. En ese ascenso, $T=D$ deja de ser exacto: la expresión logarítmica usual desprecia el empuje adicional de trayectoria. La corrección se deriva más adelante en este capítulo. No es una fórmula exacta para conservar simultáneamente altitud, TAS y $C_L$ en una atmósfera fija.

Puede usarse como aproximación de un intervalo corto a altitud fija si la variación de $L/D$ es pequeña; entonces se debe declarar el error o el criterio de aproximación. Es preferible utilizar la integral correcta si ya se dispone de una polar.

\section{Caso II: altitud y $C_L$ constantes}
Con densidad constante y $C_L$ fijo,
\begin{equation}
V(W)=\sqrt{\frac{2W}{\rho SC_L}},\qquad D(W)=W\frac{C_D}{C_L}.
\end{equation}
El tiempo conserva la forma logarítmica, pero el alcance no:
\begin{align}
E&=\frac{C_L}{g_0c_TC_D}\ln\frac{W_i}{W_f},\\
R_a&=\frac{2}{g_0c_T}\sqrt{\frac{2}{\rho S}}
\frac{\sqrt{C_L}}{C_D}(\sqrt{W_i}-\sqrt{W_f}).
\label{eq:jet-CL}
\end{align}
En términos de masa, se sustituyen $W_i=g_0m_i$ y $W_f=g_0m_f$. La velocidad final es $V_f=V_i\sqrt{m_f/m_i}$. La misión tarda más que una división por la velocidad inicial, porque el avión se ralentiza durante el consumo.

\subsection{Óptimos diferentes}
Si $c_T$ es constante y no existen restricciones adicionales, maximizar autonomía exige maximizar $C_L/C_D$:
\begin{equation}
C_{L,E}=\sqrt{C_{D0}/k}.
\end{equation}
Para maximizar alcance a densidad y $C_L$ constantes se maximiza $\sqrt{C_L}/C_D$. Al derivar,
\begin{equation}
\frac{\dd}{\dd C_L}\left(\frac{\sqrt{C_L}}{C_{D0}+kC_L^2}\right)=0
\quad\Longrightarrow\quad C_{L,R}=\sqrt{\frac{C_{D0}}{3k}}.
\end{equation}
La velocidad de máximo alcance específico es $3^{1/4}$ veces la de mínima resistencia para el mismo peso y densidad. La condición puede cambiar al incluir consumo específico dependiente de Mach o gases, límite de Mach, variación de polar o restricciones operativas.

\section{Caso III: altitud y TAS constantes}
Ahora $C_L=W/(q_\infty S)$ cambia con el peso. Conviene escribir
\begin{equation}
D(W)=A_V+B_VW^2,\quad
A_V=\tfrac12\rho V^2SC_{D0},\quad B_V=\frac{2k}{\rho V^2S}.
\end{equation}
$A_V$ es una fuerza y $B_V$ tiene unidad $\mathrm{N^{-1}}$. Para $c_T$ constante,
\begin{equation}
E=\frac{1}{g_0c_T\sqrt{A_VB_V}}
\left[\arctan\!\left(W_i\sqrt{\frac{B_V}{A_V}}\right)
-\arctan\!\left(W_f\sqrt{\frac{B_V}{A_V}}\right)\right],
\label{eq:jet-arctan}
\end{equation}
y $R_a=VE$. La raíz $W\sqrt{B_V/A_V}$ es adimensional. La derivada de la primitiva recupera $1/[g_0c_T(A_V+B_VW^2)]$; esta comprobación se ejecuta simbólicamente en el proyecto.

\subsection{Inversión para encontrar combustible}
Si se prescribe tiempo $E$, sea $\Omega=g_0c_T\sqrt{A_VB_V}E$. Entonces
\begin{equation}
W_f=\sqrt{\frac{A_V}{B_V}}
\tan\!\left[\arctan\!\left(W_i\sqrt{\frac{B_V}{A_V}}\right)-\Omega\right].
\end{equation}
La solución debe conservar $W_f$ por encima del peso seco y las reservas prescritas. Una tangente negativa o una rama que atraviese una singularidad no se interpreta como masa física. Para un alcance prescrito se usa $E=R_a/V$ sólo si TAS y ausencia de viento cumplen lo declarado.

\section{La ley de empuje constante}
Si $T=T_c$ y $c_T$ es constante, el gasto $c_TT_c$ es constante. La masa decrece linealmente:
\begin{equation}
m(t)=m_i-c_TT_ct.
\end{equation}
La velocidad debe satisfacer en cada instante $D(V,W)=T_c$. Al disminuir peso, las dos ramas de equilibrio cambian. El alcance se obtiene integrando la velocidad de la rama elegida. Mantener palanca fija no es necesariamente mantener empuje fijo; esta ley describe un modelo o un control que realmente conserve $T_c$.

El dominio termina cuando se alcanza la reserva, un límite de velocidad o una condición de trim. Un cálculo que consume todo el tanque antes del final de una primera fase no permite continuar una segunda fase propulsada con masa inferior a la seca.

\ilustra{leyes-pilotaje}{Evolución de TAS normalizada con masa para dos leyes de vuelo a densidad fija. Conservar $C_L$ exige reducir velocidad; conservar TAS exige reducir $C_L$.}

\section{Altitud óptima: no hay una respuesta universal}
En el caso de densidad y $C_L$ constantes, el alcance contiene $\rho^{-1/2}/c_T$. Si se postula $c_T=c_{T0}\sigma^y$, el factor ambiental es proporcional a $\sigma^{-(1/2+y)}$. Esta tendencia sugiere subir cuando $y>-1/2$, pero sólo hasta que empuje, Mach, techo, polar, buffet o restricciones impidan mantener la condición. La autonomía a $C_L$ constante depende de $1/c_T$, sin el mismo factor de velocidad.

Un óptimo sin restricciones puede quedar en una frontera no física. No se deduce ``siempre conviene la máxima altura'' a partir de una sola dependencia. En una misión completa el ascenso consume tiempo y combustible; el beneficio de crucero debe compensarlo.

\section{Viento y alcance terrestre}
Con viento uniforme de cola $U$ alineado con la ruta,
\begin{equation}
R_g=R_a+UE.
\end{equation}
La autonomía del modelo de consumo no cambia por una traslación uniforme de la masa de aire si se mantiene la misma ley aérea. El combustible para una distancia terrestre sí cambia. Con $C_L$ constante, $V$ varía, así que no debe sustituirse la velocidad por $V_i+U$ en toda una fórmula derivada para otra ley.

Con viento cruzado se integra el avance proyectado manteniendo la ruta deseada. Si $U_E$ es constante, la componente aérea útil vale $\sqrt{V(t)^2-U_E^2}$ mientras sea real. Una TAS decreciente puede hacer inviable la corrección de deriva al final aunque fuera posible al principio.

\section{Una misión por fases}
Cada fase recibe la masa final y posición final de la anterior. Una tabla mínima registra $m_i,m_f,t_i,t_f,h_i,h_f,R$ y modelo de consumo. Si se descarga carga de pago, se registra como evento de masa distinto del combustible. Si se vuelve sin repostar, el combustible de retorno debe estar a bordo en la salida; no se calcula el viaje de ida con una masa que omite el combustible necesario para volver.

Para dimensionar combustible con peso seco y ruta prescritos, la masa inicial es una incógnita que también aumenta la resistencia. Se resuelve una ecuación de cierre $m_i=m_{\rm seco}+m_{\rm combustible}(m_i)+m_{\rm reserva}$. Una iteración que no converge o viola capacidad de tanque indica un diseño o misión inviables bajo esas hipótesis, no un motivo para fijar masa constante.

\caso{Ejemplo original y comprobación por ODE}{Se emplea $m_i=4800\SIu{kg}$, $m_f=3900\SIu{kg}$, $S=28\SIu{m^2}$, polar $0.024+0.046C_L^2$, $h=4500\SIu m$ ISA, $V=105\SIu{m\,s^{-1}}$ y $c_T=1.65\times10^{-5}\SIu{kg/(N\,s)}$. La ecuación~\eqref{eq:jet-arctan} se compara con integrar $\dot m=-c_TD(V,m)$ y $\dot x=V$ hasta el evento $m=3900\SIu{kg}$. Se repite con pasos máximos de 300, 120 y 30 s. Las tres comparaciones pasan un error relativo de tiempo de $2\times10^{-7}$. No se compara con una ecuación logarítmica de hipótesis diferentes.}

\ilustra{jet-consumo}{Alcance calculado al variar combustible consumido: una ley conserva $C_L$ de máximo alcance y otra conserva TAS de $105\SIu{m\,s^{-1}}$. Ambas usan masa variable y la misma polar docente.}

\section{Integración explícita con empuje constante}
\subsection{Parametrizar las dos ramas}
En altitud fija, se escribe $D=AV^2+CW^2/V^2$, donde $A=\rho SC_{D0}/2$ y $C=2k/(\rho S)$. Con empuje constante $T_c$,
\begin{equation}
AV^4-T_cV^2+CW^2=0.
\end{equation}
La existencia exige $W\leq T_c/(2\sqrt{AC})$. Se define
\begin{equation}
W=K_W\sin2\zeta,\quad K_W=\frac{T_c}{2\sqrt{AC}},\qquad 0<\zeta\leq\frac\pi4.
\end{equation}
Entonces las dos velocidades son
\begin{equation}
V_+=\sqrt{T_c/A}\cos\zeta,\qquad
V_-=\sqrt{T_c/A}\sin\zeta.
\end{equation}
La rama alta aumenta velocidad al disminuir peso y $\zeta$; la baja la reduce. Ambas se encuentran en el mínimo de resistencia cuando $\zeta=\pi/4$. La selección de una rama debe conservar pérdida, Mach y trim durante todo el recorrido.

\subsection{Tiempo y alcance de cada rama}
El peso cambia linealmente: $\dot W=-g_0c_TT_c$. Por tanto,
\begin{equation}
E=\frac{W_i-W_f}{g_0c_TT_c}.
\end{equation}
La distancia se obtiene con $\int V\,\dd W/(g_0c_TT_c)$. Como $\dd W=2K_W\cos2\zeta\,\dd\zeta$, la rama alta exige
\begin{equation}
\int\cos\zeta\cos2\zeta\,\dd\zeta
=\sin\zeta-\frac23\sin^3\zeta.
\end{equation}
De ahí,
\begin{equation}
R_+=\frac{\sqrt{T_c/A}}{g_0c_T\sqrt{AC}}
\left[\sin\zeta-\frac23\sin^3\zeta\right]_{\zeta_f}^{\zeta_i}.
\end{equation}
Para la rama baja,
\begin{equation}
R_-=\frac{\sqrt{T_c/A}}{g_0c_T\sqrt{AC}}
\left[\cos\zeta-\frac23\cos^3\zeta\right]_{\zeta_f}^{\zeta_i}.
\end{equation}
Ambas primitivas se verifican por diferenciación. El tiempo puede coincidir para las dos ramas porque el gasto es idéntico, mientras la distancia cambia con velocidad. Igual consumo total no implica igual alcance.

\section{Crucero con ascenso: la corrección que suele omitirse}
La fórmula logarítmica con V y eficiencia constantes se derivó suponiendo $T=D$. Un crucero que cambia de altura no cumple esa igualdad exactamente. Para trayectoria rectilínea de ángulo constante y velocidad constante,
\begin{equation}
L=W\cos\gamma,\qquad
T=W\left(\frac{\cos\gamma}{E_a}+\sin\gamma\right).
\end{equation}
Si además $E_a,c_T,\gamma$ son constantes, el tiempo y alcance horizontal son
\begin{equation}
E=\frac{\ln(W_i/W_f)}{g_0c_T(\cos\gamma/E_a+\sin\gamma)},\qquad
R_H=V\cos\gamma\,E.
\end{equation}
Cuando $\gamma\to0$ se recupera la expresión horizontal. En un crucero con ascenso real, $\gamma$ puede variar y la densidad debe ser compatible con peso, velocidad, sustentación y atmósfera. La fórmula anterior es un caso condicionado, no una solución automática de crucero ascendente ISA.

La aproximación habitual $T\simeq D$ en ascenso muy lento omite el término $W\sin\gamma$. Su importancia relativa frente a D es aproximadamente $E_a\gamma$ para ángulos pequeños. Incluso un ángulo pequeño puede producir una corrección apreciable si la eficiencia es grande. Por ejemplo, con $E_a=18$ y $\gamma=0.2^\circ$, la escala relativa es del orden del 6\%. El error no se juzga sólo mirando que el ángulo sea menor que un grado.

\section{Óptimo local frente a óptimo de misión}
Maximizar alcance específico local consiste en maximizar $V/(c_TD)$ en el estado actual. Un óptimo integral puede coincidir con mantener un CL óptimo si las dependencias se separan y no existen restricciones que acoplen el futuro. Si el motor cambia consumo con altura, hay reservas, escalones de nivel, tramos restringidos o coste de ascenso, la decisión actual puede modificar los estados posteriores.

Una formulación de misión usa estados $m,h,x$ y controles como velocidad y gases. Se integra cada candidato conservando restricciones y se compara el objetivo total. Para una búsqueda numérica sencilla se pueden discretizar altitudes y leyes de velocidad, resolver cada misión y conservar sólo las factibles. El resultado se declara óptimo dentro de esa familia de candidatos, sin confundirlo con un óptimo global de todas las trayectorias posibles.

\subsection{Combustible por distancia terrestre}
Con viento de cola U uniforme, el alcance específico terrestre local es $(V+U)/(c_TD)$ en vuelo horizontal alineado. Su condición óptima respecto a V incluye derivadas de consumo y resistencia:
\begin{equation}
\frac1{V+U}-\frac1{c_T}\frac{\partial c_T}{\partial V}
-\frac1D\frac{\partial D}{\partial V}=0.
\end{equation}
La derivada se evalúa conservando peso y condición ambiental del punto. Con cT constante y calma se reduce a $D-VD'=0$, que para $D=AV^2+B/V^2$ conduce a $V^4=3B/A$. Esta comprobación conecta el criterio local de alcance con el CL de máximo alcance de reactor.

\section{Sensibilidad de una misión y reservas}
Una misión que llega exactamente con masa seca es sensible a cualquier variación adversa de viento o consumo. Para explorar sensibilidad se recalcula con cT aumentado, polar modificada o viento desfavorable y se registra qué restricciones se activan. Si se prescribe reserva como masa, se conserva en el peso final. Si se prescribe como tiempo de espera, se resuelve esa fase con su ley y se transporta su combustible durante las fases anteriores.

No se debe añadir un porcentaje genérico a una fórmula y presentarlo como reserva operativa universal. El análisis académico puede utilizar una reserva dada por el enunciado y estudiar cómo cambia la masa inicial. La definición operacional y normativa de reservas requiere el contexto y documentos aplicables y queda fuera de estas expresiones de diseño docente.


\section{Taller resuelto: Dimensionar combustible por cierre de misión}

\noindent\textbf{Nivel:} Intermedio.

\apartado{1. Enunciado}{Una fase de reactor debe recorrer 1400 km en calma a 7000 m ISA y CL de máximo alcance. La masa final, que incluye carga y reserva, es 6500 kg. Con S=42 m², polar 0.024+0.044CL² y cT=1.7e-5 kg/(N s), determine masa inicial y combustible de la fase.}

\apartado{2. Datos}{El combustible de la fase es desconocido y forma parte de masa inicial.}

\apartado{3. Incógnitas}{Las magnitudes y comprobaciones solicitadas en el caso.}

\apartado{4. Sistema de referencia}{Crucero horizontal y TAS variable para CL fijo.}

\apartado{5. Hipótesis físicas}{Polar, densidad y consumo constantes; no hay ascenso ni transición dentro de la fase.}

\apartado{6. Modelo matemático}{Dimensionar combustible por cierre de misión}

\apartado{7. Ecuaciones simbólicas}{Se parte de los balances y definiciones siguientes.}

\begin{equation}R=K_m(\sqrt{m_i}-\sqrt{m_f})\end{equation}

\begin{equation}m_i=(\sqrt{m_f}+R/K_m)^2\end{equation}

\begin{equation}m_{comb}=m_i-m_f\end{equation}

\apartado{8. Desarrollo algebraico}{Las operaciones se realizan en este orden, conservando las condiciones del modelo.}

\begin{enumerate}

\item Se calcula CL de máximo alcance y el coeficiente Km de la ley.

\item Se despeja masa inicial antes de restar combustible; no se utiliza una masa media arbitraria.

\item Se verifican velocidades extremas y que el alcance se reconstruya con la masa hallada.

\end{enumerate}

\apartado{9. Sustitución numérica}{Se conservan los resultados intermedios sin redondear y se muestran todas las operaciones escalares disponibles.}

\textbf{Desarrollo de las magnitudes intermedias.} De $dm/dt=-c_TD=-c_Tmg_0/(L/D)$, separar variables da $E=(L/D)\ln(m_i/m_f)/(g_0c_T)$. Con CL fijo, $V=\sqrt{2g_0m/(\rho SC_L)}$; sustituir en $dx=Vdt$ deja integrar $m^{-1/2}$, cuya primitiva es $2\sqrt m$. Por eso $R=K_m(\sqrt{m_i}-\sqrt{m_f})$, con $K_m=2(L/D)\sqrt{2g_0/(\rho SC_L)}/(g_0c_T)$. Para reactor a CL constante, máximo alcance se obtiene maximizando $\sqrt{C_L}/C_D$, que da $C_L=\sqrt{C_{D0}/(3k)}$. Despejar la masa inicial requiere primero $\sqrt{m_i}=\sqrt{m_f}+R/K_m$ y después elevar al cuadrado.\begin{enumerate}

\item Elegir el óptimo de alcance de reactor.

\[\fitmath{C_L=\sqrt{\frac{0.024}{3\cdot 0.044}}\approx 0.4264014\,}\]

\item Evaluar resistencia adimensional.

\[\fitmath{C_D=0.024+0.044\cdot {C_L}^{2}}\]

\[\fitmath{C_D=0.024+0.044\cdot {(0.4264014327)}^{2}\approx 0.032\,}\]

\item Obtener eficiencia de esa condición.

\[\fitmath{\mathrm{LD}=\frac{C_L}{C_D}}\]

\[\fitmath{\mathrm{LD}=\frac{(0.4264014327)}{(0.032)}\approx 13.32504\,}\]

\item Calcular el coeficiente completo de la integral.

\[\fitmath{\mathrm{Km}=\frac{2\cdot \mathrm{LD}}{9.80665\cdot 1.7\times10^{-5}}\cdot \sqrt{\frac{2\cdot 9.80665}{0.589500721\cdot 42\cdot C_L}}}\]

\[\fitmath{\mathrm{Km}=\frac{2\cdot (13.32504477)}{9.80665\cdot 1.7\times10^{-5}}\cdot \sqrt{\frac{2\cdot 9.80665}{0.589500721\cdot 42\cdot (0.4264014327)}}\approx 217885.5\,\mathrm{m/kg^{.5}}}\]

\item Cerrar el alcance prescrito antes de estimar combustible.

\[\fitmath{\mathrm{mi}={\left(\sqrt{6500}+\frac{1400000}{\mathrm{Km}}\right)}^{2}}\]

\[\fitmath{\mathrm{mi}={\left(\sqrt{6500}+\frac{1400000}{(217885.4692)}\right)}^{2}\approx 7577.349\,\mathrm{kg}}\]

\item Restar masa final que ya incluye reserva.

\[\fitmath{\mathrm{fuel}=\mathrm{mi}-6500}\]

\[\fitmath{\mathrm{fuel}=(7577.34935)-6500\approx 1077.349\,\mathrm{kg}}\]

\item Comprobar velocidad inicial.

\[\fitmath{\mathrm{Vi}=\sqrt{\frac{2\cdot \mathrm{mi}\cdot 9.80665}{0.589500721\cdot 42\cdot C_L}}}\]

\[\fitmath{\mathrm{Vi}=\sqrt{\frac{2\cdot (7577.34935)\cdot 9.80665}{0.589500721\cdot 42\cdot (0.4264014327)}}\approx 118.6473\,\mathrm{m/s}}\]

\item Comprobar velocidad final.

\[\fitmath{\mathrm{Vf}=\sqrt{\frac{2\cdot 6500\cdot 9.80665}{0.589500721\cdot 42\cdot C_L}}}\]

\[\fitmath{\mathrm{Vf}=\sqrt{\frac{2\cdot 6500\cdot 9.80665}{0.589500721\cdot 42\cdot (0.4264014327)}}\approx 109.8894\,\mathrm{m/s}}\]

\end{enumerate}





\apartado{10. Resultados con unidades}{Los resultados pertenecen al modelo y al intervalo de datos declarados.}

\begin{center}\begin{tabular}{lr}\toprule Magnitud&Resultado\\\midrule

Masa inicial & $7577.349\,\mathrm{kg}$\\

Combustible de fase & $1077.349\,\mathrm{kg}$\\

TAS inicial & $118.6473\,\mathrm{m\,s^{-1}}$\\

TAS final & $109.8894\,\mathrm{m\,s^{-1}}$\\

\bottomrule\end{tabular}\end{center}

\apartado{11. Interpretación}{La masa final ya contiene la reserva; el combustible calculado es sólo el quemado en esta fase.}

\apartado{12. Comprobación independiente y alcance}{Resolver por raíz numérica la diferencia entre alcance calculado y prescrito; comprobar cierre de masa.}

\apartado{13. Limitaciones}{La factibilidad necesita capacidad de tanque, masa máxima, empuje y límites de Mach, que no se han dado.}

\par\smallskip\noindent{\footnotesize\textit{Datos y solución originales. Las sustituciones registradas se reproducen con SymPy y Studium; eso no establece validación experimental ni revisión experta.}}


\repaso{¿Es correcto usar el logaritmo de Breguet con la velocidad inicial para un crucero a altitud y $C_L$ constantes?}{La fórmula temporal logarítmica sí corresponde a esas hipótesis si $c_T$ y polar son constantes. El alcance exige integrar $V\propto\sqrt W$ y contiene una diferencia de raíces. Usar $V_i$ constante sobreestima el alcance de esa ley.}
\resumen{La ley de vuelo determina qué integral se resuelve. A altitud fija, $C_L$ constante conduce a raíces en alcance; TAS constante conduce a arcotangentes. Alcance y autonomía tienen óptimos distintos para reactores, y una misión exige actualizar masa y restricciones en cada fase.}
\fuentes{Guía ULE, actuaciones integrales de aerorreactores. Las integrales se derivan explícitamente desde gasto másico y polar, y se contrastan mediante SymPy y SciPy. Los textos especializados citados en la bibliografía recomendada no se presentan como secciones consultadas para estas deducciones.}
